\setcounter{chapter}{5}
\setcounter{section}{0}

\chapter{Esiti}
    Le seguenti sezioni illustrano come il Corso di Studio in Scienze e Tecnologie Multimediali possa creare una figura tecnica professionale e di come possa preparare ai vari dispositivi multimediali e informatici che si possono trovare in un Polo Universitario. \newline
    Viene poi dedicata una sezione ai possibili sbocchi lavorativi nelle P.A., nello specifico nella Scuola Secondaria di Secondo Grado e nelle Università. 
% ---------------------------------------------------------------------------------------
    \section{Preparazione del Corso di Studio}
        Avendo ottenuto il Diploma di Informatica e Telecomunicazioni, non sono un estraneo al mondo dell'informatica; nonostante tale formazione di base ho potuto personalmente esprimere come l'offerta formativa del CdS sia molto buona, fornendo un'efficace preparazione a tutte quelle casistiche al limite tra l'informatica e il multimediale che ho riscontrato durante il periodo di tirocinio presso il CEGO. \newline 
        Ho notato maggior focus da parte del CdS nelle materie tecniche e laboratoriali, soprattutto sull'insegnamento Laboratorio di Programmazione e Sistemi Multimediali sostenuto con il Docente e mio Relatore Prof. Niki Martinel in quanto tale percorso mi ha dato una preparazione superiore rispetto ad altri insegnamenti, sempre di natura tecnica. \newline
        Poiché l'Insegnamento tratta sia di programmazione e sviluppo di applicativi sia che di sistemi multimediali, veniva a essere toccato ogni aspetto tecnico che si possa immaginare: la parte riguardante i sistemi multimediali ha illustrato con cura il punto di vista storico-tecnologico dell'informatica, mentre la seconda parte ha fatto riferimento  più su strutture dati e algoritmi da sviluppare in un linguaggio ad alto livello, Python. Questo approccio multiplo della materia mi ha aiutato a capire meglio come funzionano molti applicativi e soprattutto mi ha aiutato a capire le strutture dei vari file, degli algoritmi e la manipolazione di dati e formati multimediali e come questi dati vengono gestiti dal software; durante il tirocinio mi ha aiutato a identificare meglio alcuni problemi con cui mi sono dovuto misurare. \newline
        I testi di studio inoltre mi hanno guidato non solo nella preparazione all'esame ma anche durante il periodo di tirocinio, illustrandomi le diverse tecnologie e i diversi meccanismi su cui non solo si fonda l'informatica di oggi ma che si trovano ormai ovunque.
        \par
            Durante il periodo di tirocinio ho avuto modo di conoscere diversi aspetti della vita universitaria, prima sconosciuti: ho vissuto di prima mano cosa vuol dire lavorare nell'Amministrazione di un Polo Universitario e nello specifico ho potuto cimentarmi nelle mansioni di un Tecnico Informatico. \newline
            Questo tirocinio mi ha aiutato ad ampliare le mie conoscenze tecniche e mi ha permesso di migliorare quei aspetti professionali che prima non avevo approfondito, data la natura scolastica del mio percorso. I ritmi che in alcuni momenti definirei frenetici degli interventi mi hanno dato una prospettiva nuova, prima sconosciuta, delle possibilità lavorative post-laurea. Mi ha anche permesso di ampliare i miei orizzonti per quanto riguarda altri possibili sbocchi lavorativi nel pubblico, piuttosto che nel privato. 
% ---------------------------------------------------------------------------------------
    \section{Inserimento Professionale}
        \par
            Gli inserimenti professionali possibili sono molteplici e possono variare di molto in base a ciò in cui mi poteri specializzare. Inizialmente, l'obiettivo era quello di proseguire immediatamente gli studi attraverso il conseguimento di una Laurea Magistrale ma, dopo il periodo di Tirocinio, mi sono ritrovato di fronte ad un bivio: proseguire gli studi e rimandare l'ingresso nel mondo del lavoro oppure dare in un certo senso prosecuzione all'attività di tirocinio, puntando ad un posto di lavoro come tecnico informatico. 
        \par
            Tutt'ora il proseguire gli studi costituisce una forte spinta motivante a trovare un posto di lavoro che mi permetta di conciliare entrambe le possibilità: da qui nasce il desiderio di lavorare nelle Pubbliche Amministrazioni, attraverso la partecipazione a diversi bandi sia nell'ambito delle Scuole Secondarie di Secondo Grado o magari nell'ambito della stessa Università di Udine.
        \par
            Per quanto riguarda il profilo di Tecnico Informatico presso l'Università, i bandi vengono pubblicati in una sezione apposita presente nel sito web, dove vengono catalogati e raggruppati per Dipartimento richiedente; nello specifico, l'ente universitario \cite{DARU}.
        \par
            Per le Scuole Secondarie di Secondo Grado, i posti di lavoro disponibili per un neo-laureato triennale sono principalmente due: Personale ATA di 3a Fascia e Insegnate Tecnico Pratico per la cattedra B016. \newline
            Per il profilo ATA di Assistente Tecnico il requisito di base è ancora il solo Diploma di Scuola Secondaria di Secondo Grado corrispondente allo specifico settore professionale unito alla Certificazione Internazionale di Alfabetizzazione digitale, come previsto dal \cite{CCNL-Istruzione/Ricerca}, entrato in vigore 1 maggio 2024. Tuttavia la Laurea Triennale è un titolo valutabile utile ad ottenere un punteggio aggiuntivo. Per il Personale Tecnico Amministrativo (PTA) in ambito universitario si fa riferimento al D.Lgs. 165/2001, CCNL Funzioni Centrali/Istruzione e Ricerca e al regolamento interno dell'Università.\newline
            La Laurea Triennale per gli Insegnanti Tecnico-Pratici (ITP) è stata invece stabilita dal \cite{D.L.59/2017} tuttavia la sua entrata in vigore è stata prorogata, spostando il termine ultimo dall'iniziale 2025 al 31 dicembre 2026 con il \cite{DecretoMilleproroghe}. A partire da tale termine, serviranno la Laurea Triennale oltre ai 24 CFU/CFA per accedere ai concorsi, mentre finora è stato sufficiente il solo Diploma di Scuola Secondaria di Secondo Grado. 
        \par
            Il percorso appena sopra delineato può sembrare accidentato e complesso ma risulta essere altamente poiché può dar luogo a una figura professionalmente qualificata, permettendo quasi di "chiudere un cerchio" andando virtualmente a congiungere le materie oggetto di studio all'ambiente in cui questo si è svolto, come può esserlo la scuola o l'Università. \newline
            Il percorso lavorativo risulta invece essere indubbiamente diverso nell'ambito privato, in quanto le aziende tendono a cercare figure specializzate, ad esempio sistemisti o sviluppatori, oppure Tecnici manutentori o \textit{Information Technology} (IT), che operino molto più spesso di quanto accada nell'ambito dei servizi pubblici nelle fasce orarie in cui l'azienda sia vuota allo scopo di non interromperne la produttività e diminuendo in tal modo fortemente l'interazione interpersonale con il capitale umano. 